\documentclass[11pt,a4paper]{article}

\usepackage[utf8]{inputenc}
\usepackage[T1]{fontenc}
\usepackage[margin=2.5cm]{geometry}
\usepackage{hyperref}
\usepackage{xcolor}
\usepackage{listings}
\usepackage{enumitem}
\usepackage{booktabs}
\usepackage{longtable}
\usepackage{amsmath,amssymb}
\usepackage{fancyhdr}
\usepackage{titlesec}

\definecolor{isls-blue}{HTML}{1a5276}
\definecolor{isls-orange}{HTML}{e67e22}
\definecolor{isls-green}{HTML}{27ae60}
\definecolor{isls-red}{HTML}{c0392b}
\definecolor{isls-gray}{HTML}{7f8c8d}
\definecolor{code-bg}{HTML}{f7f9fb}

\lstdefinestyle{rust}{
  language=C,
  basicstyle=\ttfamily\small,
  keywordstyle=\color{isls-blue}\bfseries,
  commentstyle=\color{isls-gray}\itshape,
  stringstyle=\color{isls-green},
  backgroundcolor=\color{code-bg},
  frame=single,
  rulecolor=\color{isls-gray!30},
  breaklines=true,
  showstringspaces=false,
  tabsize=4,
  morekeywords={fn,let,mut,pub,struct,impl,enum,use,mod,self,Self,
    crate,super,where,trait,for,in,if,else,match,return,async,await,
    Ok,Err,Some,None,Vec,String,HashMap,BTreeMap,BTreeSet,
    Result,Option,PathBuf,
    f64,f32,usize,bool,u32,u64,i64,Box,dyn,Arc},
}
\lstdefinestyle{bash}{
  language=bash,
  basicstyle=\ttfamily\small,
  backgroundcolor=\color{code-bg},
  frame=single,
  rulecolor=\color{isls-gray!30},
  breaklines=true,
  showstringspaces=false,
}

\pagestyle{fancy}
\fancyhf{}
\fancyhead[L]{\small\textcolor{isls-gray}{ISLS I4 Spec --- Cross-Language Barbara}}
\fancyhead[R]{\small\textcolor{isls-gray}{v1.0.0 --- \today}}
\fancyfoot[C]{\thepage}

\titleformat{\section}
  {\Large\bfseries\color{isls-blue}}{\thesection}{1em}{}
\titleformat{\subsection}
  {\large\bfseries\color{isls-blue!80}}{\thesubsection}{1em}{}
\titleformat{\subsubsection}
  {\normalsize\bfseries\color{isls-blue!60}}{\thesubsubsection}{1em}{}

\hypersetup{
  colorlinks=true,
  linkcolor=isls-blue,
  urlcolor=isls-orange,
}

\begin{document}

\begin{titlepage}
\centering
\vspace*{2.5cm}

{\Huge\bfseries\color{isls-blue} ISLS I4 Specification}\\[0.5cm]
{\LARGE\color{isls-orange} Cross-Language Barbara}\\[0.3cm]
{\large\color{isls-gray} Multi-Language Parsing, Sprachunabh\"angige\\
Resonites, Superpositions-Verst\"arkung}\\[2cm]

{\large
\begin{tabular}{rl}
\textbf{System:}       & ISLS --- Intelligent Semantic Ledger Substrate \\
\textbf{Version:}      & v6.3 (I4) \\
\textbf{Author:}       & Sebastian Klemm \\
\textbf{Date:}         & \today \\
\textbf{Status:}       & Specification \\
\textbf{Depends:}      & I3 (complete) \\
\textbf{Branch:}       & \texttt{claude/implement-isls-i4-spec} \\
\textbf{Crate:}        & \texttt{isls-reader} (modified) \\
\textbf{Agent:}        & Sonnet 4.6 (klar abgegrenzte Parser-Arbeit) \\
\end{tabular}
}

\vfill

{\small\color{isls-gray}
Mithras Substructure University\\
Repository: \texttt{https://github.com/LashSesh/graphity}\\
License: MIT
}
\end{titlepage}

\tableofcontents
\newpage


% ══════════════════════════════════════════════════════════════════════════════
\section{Executive Summary}
% ══════════════════════════════════════════════════════════════════════════════

Barbara (\texttt{isls-reader}) parst aktuell nur Rust. Die 
\"Olbohrinsel scrapt nur Rust-Repos von GitHub. Das schr\"ankt 
den Norm-Korpus auf $\sim$5\% des GitHub-\"Okosystems ein.

I4 erweitert Barbara um drei weitere Sprachen: \textbf{Python}, 
\textbf{TypeScript} und \textbf{Go}. Zusammen decken Rust + 
Python + TypeScript + Go $\sim$60\% aller relevanten 
Software-Repositories auf GitHub ab.

Der Schl\"ussel: die Resonites sind sprachunabh\"angig. 
Eine Funktion \texttt{get\_connection(pool)} produziert 
denselben Resonite ob sie in Rust, Python, TypeScript oder 
Go geschrieben ist. Die Normen die aus cross-language 
Scraping kristallisieren sind robuster, weil sprachspezifische 
Fehler sich nicht korreliert fortpflanzen (Theorem 4.1 der 
Metamorphose-Formalisierung).

\subsection{Was I4 Produziert}

\begin{enumerate}[nosep]
\item \textbf{Python-Parser:} Extrahiert Functions, Classes, 
      Imports, Decorators aus \texttt{.py}-Dateien.
\item \textbf{TypeScript-Parser:} Extrahiert Functions, 
      Interfaces, Classes, Imports aus \texttt{.ts/.tsx}-Dateien.
\item \textbf{Go-Parser:} Extrahiert Functions, Structs, 
      Interfaces, Imports aus \texttt{.go}-Dateien.
\item \textbf{Unified Resonites:} Alle vier Parser produzieren 
      dieselben \texttt{Resonite}-Varianten. Keine 
      sprachspezifischen Resonite-Typen.
\item \textbf{Language-Aware Scraping:} Die \"Olbohrinsel 
      scrapt ohne \texttt{language:rust} Filter und verdaut 
      Repos in jeder der vier Sprachen.
\item \textbf{Cross-Language Norm-Tag:} Normen die aus 
      2+ Sprachen kristallisieren werden als 
      \texttt{cross\_language: true} markiert.
\end{enumerate}


% ══════════════════════════════════════════════════════════════════════════════
\section{Architektur: Sprachunabh\"angige Resonites}
\label{sec:arch}
% ══════════════════════════════════════════════════════════════════════════════

\subsection{Das Prinzip}

Resonites abstrahieren \"uber Syntax. Sie beschreiben 
\emph{was} existiert (eine Funktion, ein Typ, ein Import), 
nicht \emph{wie} es geschrieben ist. Diese Abstraktion ist 
der Grund warum cross-language Normen m\"oglich sind.

\subsection{Beispiel}

Dasselbe Pattern in vier Sprachen:

\begin{lstlisting}[style=bash]
# Rust:
pub async fn get_user(pool: &PgPool, id: Uuid) -> Result<User>

# Python:
async def get_user(pool: AsyncSession, id: UUID) -> User

# TypeScript:
export async function getUser(pool: Pool, id: string): Promise<User>

# Go:
func GetUser(pool *pgxpool.Pool, id uuid.UUID) (*User, error)
\end{lstlisting}

Alle vier produzieren denselben Resonite:

\begin{lstlisting}[style=rust]
Resonite::Fn {
    name: "get_user",    // normalisiert zu snake_case
    arity: 2,
    return_type: "User", // nur der Kerntyp
    is_pub: true,
    is_async: true,
}
\end{lstlisting}

\subsection{Normalisierungsregeln}

\begin{longtable}{lp{8cm}}
\toprule
\textbf{Aspekt} & \textbf{Normalisierung} \\
\midrule
Funktionsnamen 
  & \texttt{camelCase} $\to$ \texttt{snake\_case}. 
    \texttt{GetUser} $\to$ \texttt{get\_user}. 
    \texttt{getUser} $\to$ \texttt{get\_user}. \\
Typnamen
  & Beibehalten in PascalCase. \texttt{User} bleibt 
    \texttt{User}. \\
Return-Typen
  & Wrapper entfernen: \texttt{Result<User>} $\to$ 
    \texttt{User}. \texttt{Promise<User>} $\to$ 
    \texttt{User}. \texttt{(*User, error)} $\to$ 
    \texttt{User}. \\
Sichtbarkeit
  & \texttt{pub}/\texttt{export}/Gro{\ss}buchstabe (Go) 
    $\to$ \texttt{is\_pub: true}. \\
Async
  & \texttt{async fn}/\texttt{async def}/\texttt{async function} 
    $\to$ \texttt{is\_async: true}. Go hat kein async 
    Keyword $\to$ \texttt{is\_async: false} (Goroutines 
    sind implizit). \\
Imports
  & Pfad beibehalten. \texttt{is\_external}: 
    stdlib $=$ false, alles andere $=$ true. \\
\bottomrule
\end{longtable}


% ══════════════════════════════════════════════════════════════════════════════
\section{W1: Python-Parser}
\label{sec:python}
% ══════════════════════════════════════════════════════════════════════════════

\subsection{Was extrahiert wird}

\begin{longtable}{lp{4cm}p{5cm}}
\toprule
\textbf{Python-Konstrukt} & \textbf{Resonite} & \textbf{Mapping} \\
\midrule
\texttt{def foo(a, b):}
  & \texttt{Resonite::Fn}
  & name, arity=2, is\_async=false \\
\texttt{async def foo(a):}
  & \texttt{Resonite::Fn}
  & is\_async=true \\
\texttt{class User:}
  & \texttt{Resonite::Type}
  & kind=Struct, field\_count from \_\_init\_\_ \\
\texttt{class Handler(ABC):}
  & \texttt{Resonite::Type}
  & kind=Trait (abstract base class) \\
\texttt{from x import y}
  & \texttt{Resonite::Import}
  & path=x.y, is\_external from known stdlibs \\
\texttt{import os}
  & \texttt{Resonite::Import}
  & path=os, is\_external=false (stdlib) \\
\texttt{@dataclass}
  & \texttt{Resonite::Type}
  & derive\_count += 1 \\
\texttt{@app.route}
  & (ignored)
  & Decorator-Routing not a structural element \\
\bottomrule
\end{longtable}

\subsection{Parsing-Strategie}

Kein vollst\"andiger Python-AST-Parser. Ein \textbf{Regex-basierter 
Lightweight-Parser} der die oben genannten Konstrukte 
erkennt. Python hat eine konsistente Syntax f\"ur Definitionen:

\begin{lstlisting}[style=rust]
// Patterns (Regex)
let fn_pattern = r"^\s*(async\s+)?def\s+(\w+)\s*\(([^)]*)\)";
let class_pattern = r"^\s*class\s+(\w+)(?:\(([^)]*)\))?:";
let import_from = r"^\s*from\s+([\w.]+)\s+import\s+(.+)";
let import_direct = r"^\s*import\s+([\w.]+)";
let decorator = r"^\s*@(\w+)";
\end{lstlisting}

Zeile f\"ur Zeile. Kein Kontext-Tracking (keine Verschachtelung, 
keine Scopes). Das reicht f\"ur Resonite-Extraktion weil 
Resonites nur Top-Level-Deklarationen sind.

\subsection{Field-Count f\"ur Classes}

Python-Klassen haben Felder in \texttt{\_\_init\_\_}:

\begin{lstlisting}[style=bash]
class User:
    def __init__(self, name: str, email: str, age: int):
        self.name = name
        self.email = email
        self.age = age
\end{lstlisting}

field\_count = Anzahl der \texttt{self.X = X} Zuweisungen 
in \texttt{\_\_init\_\_}. Alternativ: arity von 
\texttt{\_\_init\_\_} minus 1 (self).

F\"ur \texttt{@dataclass}: field\_count = Anzahl der 
annotierten Klassenattribute.

\subsection{Bekannte Standard-Bibliotheken (is\_external)}

\begin{lstlisting}[style=rust]
const PYTHON_STDLIB: &[&str] = &[
    "os", "sys", "re", "json", "math", "datetime",
    "collections", "itertools", "functools", "typing",
    "pathlib", "io", "abc", "enum", "dataclasses",
    "asyncio", "logging", "unittest", "hashlib",
    "urllib", "http", "socket", "threading",
    "multiprocessing", "subprocess", "shutil",
    "copy", "pickle", "csv", "xml", "html",
    "argparse", "configparser", "contextlib",
];
\end{lstlisting}


% ══════════════════════════════════════════════════════════════════════════════
\section{W2: TypeScript-Parser}
\label{sec:typescript}
% ══════════════════════════════════════════════════════════════════════════════

\subsection{Was extrahiert wird}

\begin{longtable}{lp{4cm}p{5cm}}
\toprule
\textbf{TS-Konstrukt} & \textbf{Resonite} & \textbf{Mapping} \\
\midrule
\texttt{function foo(a, b)}
  & \texttt{Resonite::Fn}
  & name (snake\_case), arity=2 \\
\texttt{export async function}
  & \texttt{Resonite::Fn}
  & is\_pub=true, is\_async=true \\
\texttt{const foo = (a) =>}
  & \texttt{Resonite::Fn}
  & Arrow function, arity=1 \\
\texttt{interface User \{}
  & \texttt{Resonite::Type}
  & kind=Trait, field\_count from body \\
\texttt{class UserService \{}
  & \texttt{Resonite::Type}
  & kind=Struct \\
\texttt{type Status = ...}
  & \texttt{Resonite::Type}
  & kind=Enum (if union) \\
\texttt{import \{ x \} from 'y'}
  & \texttt{Resonite::Import}
  & path=y.x \\
\texttt{import x from 'y'}
  & \texttt{Resonite::Import}
  & path=y \\
\bottomrule
\end{longtable}

\subsection{Parsing-Strategie}

Regex-basiert wie Python:

\begin{lstlisting}[style=rust]
let fn_decl = r"(?:export\s+)?(?:async\s+)?function\s+(\w+)\s*\(([^)]*)\)";
let arrow_fn = r"(?:export\s+)?(?:const|let)\s+(\w+)\s*=\s*(?:async\s+)?\(([^)]*)\)\s*=>";
let interface_decl = r"(?:export\s+)?interface\s+(\w+)";
let class_decl = r"(?:export\s+)?class\s+(\w+)";
let type_decl = r"(?:export\s+)?type\s+(\w+)\s*=";
let import_named = r"import\s+\{([^}]+)\}\s+from\s+['\"]([^'\"]+)['\"]";
let import_default = r"import\s+(\w+)\s+from\s+['\"]([^'\"]+)['\"]";
\end{lstlisting}

\subsection{Field-Count}

F\"ur Interfaces und Classes: z\"ahle Zeilen innerhalb 
der geschweiften Klammern die ein Feld-Pattern haben:

\begin{lstlisting}[style=rust]
let field_pattern = r"^\s+(\w+)\s*[?:]";
\end{lstlisting}

\subsection{Bekannte Packages (is\_external)}

In TypeScript: Imports mit relativen Pfaden 
(\texttt{./}, \texttt{../}) sind intern. Alles andere 
(npm Packages) ist extern. Node-Builtins (\texttt{fs}, 
\texttt{path}, \texttt{http}, etc.) sind stdlib.

\begin{lstlisting}[style=rust]
fn is_external_ts(path: &str) -> bool {
    !path.starts_with('.') && !path.starts_with('/')
}

const TS_STDLIB: &[&str] = &[
    "fs", "path", "http", "https", "url", "crypto",
    "os", "child_process", "stream", "buffer",
    "events", "util", "assert", "net", "tls",
    "querystring", "readline", "zlib",
];
\end{lstlisting}


% ══════════════════════════════════════════════════════════════════════════════
\section{W3: Go-Parser}
\label{sec:go}
% ══════════════════════════════════════════════════════════════════════════════

\subsection{Was extrahiert wird}

\begin{longtable}{lp{4cm}p{5cm}}
\toprule
\textbf{Go-Konstrukt} & \textbf{Resonite} & \textbf{Mapping} \\
\midrule
\texttt{func Foo(a, b int)}
  & \texttt{Resonite::Fn}
  & is\_pub = Gro{\ss}buchstabe \\
\texttt{func (s *Svc) Foo()}
  & \texttt{Resonite::Fn}
  & Method-Receiver ignoriert, arity=0 \\
\texttt{type User struct \{}
  & \texttt{Resonite::Type}
  & kind=Struct \\
\texttt{type Handler interface \{}
  & \texttt{Resonite::Type}
  & kind=Trait \\
\texttt{import "fmt"}
  & \texttt{Resonite::Import}
  & path=fmt, is\_external=false \\
\texttt{import "github.com/x"}
  & \texttt{Resonite::Import}
  & is\_external=true \\
\bottomrule
\end{longtable}

\subsection{Parsing-Strategie}

\begin{lstlisting}[style=rust]
let func_decl = r"^func\s+(?:\([^)]+\)\s+)?(\w+)\s*\(([^)]*)\)";
let type_struct = r"^type\s+(\w+)\s+struct\s*\{";
let type_interface = r"^type\s+(\w+)\s+interface\s*\{";
let import_single = r"^\s*import\s+\"([^\"]+)\"";
let import_block_item = r"^\s*\"([^\"]+)\"";
\end{lstlisting}

\subsection{Sichtbarkeit}

Go hat keine Keywords f\"ur Sichtbarkeit. Gro{\ss}buchstabe 
am Anfang = public:

\begin{lstlisting}[style=rust]
fn is_pub_go(name: &str) -> bool {
    name.chars().next()
        .map(|c| c.is_uppercase())
        .unwrap_or(false)
}
\end{lstlisting}

\subsection{Go Standard-Bibliothek}

\begin{lstlisting}[style=rust]
fn is_external_go(path: &str) -> bool {
    path.contains('.') // "github.com/..." ist extern
                       // "fmt", "os", "net" sind stdlib
}
\end{lstlisting}


% ══════════════════════════════════════════════════════════════════════════════
\section{W4: Integration in Barbara}
\label{sec:integration}
% ══════════════════════════════════════════════════════════════════════════════

\subsection{Language-Enum Erweiterung}

\begin{lstlisting}[style=rust]
#[derive(Debug, Clone, Copy, PartialEq, Eq, Hash, 
         Serialize, Deserialize)]
pub enum Language {
    Rust,
    Python,
    TypeScript,
    Go,
    Unknown,
}

impl Language {
    pub fn from_extension(ext: &str) -> Self {
        match ext {
            "rs" => Language::Rust,
            "py" => Language::Python,
            "ts" | "tsx" => Language::TypeScript,
            "go" => Language::Go,
            _ => Language::Unknown,
        }
    }
    
    pub fn from_filename(name: &str) -> Self {
        if let Some(ext) = name.rsplit('.').next() {
            Self::from_extension(ext)
        } else {
            Language::Unknown
        }
    }
}
\end{lstlisting}

\subsection{Unified parse\_string()}

\begin{lstlisting}[style=rust]
/// Parse source code in any supported language.
/// Returns the same ParseResult regardless of language.
pub fn parse_string(code: &str, lang: Language) -> ParseResult {
    match lang {
        Language::Rust => parse_rust(code),
        Language::Python => parse_python(code),
        Language::TypeScript => parse_typescript(code),
        Language::Go => parse_go(code),
        Language::Unknown => ParseResult::empty(),
    }
}
\end{lstlisting}

\subsection{Unified parse\_directory()}

\begin{lstlisting}[style=rust]
/// Parse all supported files in a directory tree.
pub fn parse_directory(path: &Path) -> DirectoryResult {
    let mut result = DirectoryResult::new();
    
    for entry in WalkDir::new(path)
        .into_iter()
        .filter_entry(|e| !is_hidden(e) && !is_vendor(e))
    {
        let entry = match entry {
            Ok(e) => e,
            Err(_) => continue,
        };
        
        if !entry.file_type().is_file() { continue; }
        
        let lang = Language::from_filename(
            entry.file_name().to_str().unwrap_or("")
        );
        
        if lang == Language::Unknown { continue; }
        
        let code = match std::fs::read_to_string(entry.path()) {
            Ok(c) => c,
            Err(_) => continue,
        };
        
        let parsed = parse_string(&code, lang);
        result.add_file(entry.path(), lang, parsed);
    }
    
    result
}
\end{lstlisting}

\subsection{Ignorierte Verzeichnisse}

\begin{lstlisting}[style=rust]
fn is_vendor(entry: &DirEntry) -> bool {
    let name = entry.file_name().to_str().unwrap_or("");
    matches!(name, 
        "node_modules" | "vendor" | "venv" | ".venv" |
        "__pycache__" | ".git" | "target" | "dist" |
        "build" | ".tox" | ".mypy_cache" | ".pytest_cache" |
        "pkg" | "bin" | "testdata"
    )
}
\end{lstlisting}

\subsection{Resonite-Normalisierung}

Alle Parser produzieren \texttt{Resonite}-Varianten mit 
normalisierten Namen:

\begin{lstlisting}[style=rust]
/// Normalize a function/method name to snake_case.
pub fn normalize_name(name: &str) -> String {
    // camelCase -> snake_case
    // PascalCase -> snake_case  
    // already_snake -> already_snake
    let mut result = String::new();
    for (i, ch) in name.chars().enumerate() {
        if ch.is_uppercase() && i > 0 {
            let prev = name.chars().nth(i - 1);
            if prev.map(|c| c.is_lowercase()).unwrap_or(false) {
                result.push('_');
            }
        }
        result.push(ch.to_lowercase().next().unwrap_or(ch));
    }
    result
}
\end{lstlisting}


% ══════════════════════════════════════════════════════════════════════════════
\section{W5: Scraping-Pipeline Erweiterung}
\label{sec:scraping}
% ══════════════════════════════════════════════════════════════════════════════

\subsection{GitHub-Suche ohne Sprachfilter}

Aktuell sucht \texttt{discover.rs} mit 
\texttt{language:rust}. I4 entfernt diesen Filter und 
l\"asst alle Sprachen zu:

\begin{lstlisting}[style=bash]
# Vorher:
GET /search/repositories?q=event+bus+language:rust

# Nachher:
GET /search/repositories?q=event+bus+language:rust
    +language:python+language:typescript+language:go
\end{lstlisting}

Alternativ: vier separate Suchen pro Keyword (eine pro 
Sprache), um die besten Repos jeder Sprache zu finden.

\subsection{Language-Aware Topologie}

\texttt{compute\_code\_topology()} in \texttt{isls-code-topo} 
erh\"alt ein neues Feld:

\begin{lstlisting}[style=rust]
pub struct CodeTopology {
    // ... bestehende Felder ...
    pub languages: HashMap<Language, usize>,  // lang -> file count
    pub primary_language: Language,
}
\end{lstlisting}

\subsection{Cross-Language Observation}

In \texttt{observe\_and\_learn()}: die Observation speichert 
die Sprache des gescrapten Repos. Wenn ein Candidate 
Observations aus 2+ Sprachen hat, wird er als 
\texttt{cross\_language: true} markiert.

\begin{lstlisting}[style=rust]
pub struct NormCandidate {
    // ... bestehende Felder ...
    pub languages: BTreeSet<Language>,
    pub cross_language: bool,  // |languages| >= 2
}
\end{lstlisting}

\subsection{Cross-Language Norm-Tag}

Wenn ein Candidate zu einer Auto-Norm promoviert wird und 
\texttt{cross\_language == true}, bekommt die Norm ein Tag:

\begin{lstlisting}[style=rust]
pub struct Norm {
    // ... bestehende Felder ...
    pub cross_language: bool,
    pub source_languages: Vec<Language>,
}
\end{lstlisting}

Im CLI und Studio angezeigt:

\begin{lstlisting}[style=bash]
ISLS-NORM-AUTO-0005  "ConnectionPool"  auto  Molecule
    cross-language: Rust, Python, Go
    domains: sqlx, asyncpg, pgxpool
    observations: 12, consistency: 1.00
\end{lstlisting}


% ══════════════════════════════════════════════════════════════════════════════
\section{Neue Dateien}
% ══════════════════════════════════════════════════════════════════════════════

\begin{longtable}{lp{10cm}}
\toprule
\textbf{File} & \textbf{Purpose} \\
\midrule
\texttt{isls-reader/src/python.rs}
  & Python-Parser: \texttt{parse\_python()}, Regex-Patterns, 
    Stdlib-Liste. \\
\texttt{isls-reader/src/typescript.rs}
  & TypeScript-Parser: \texttt{parse\_typescript()}, 
    Arrow-Functions, Interface/Class/Type-Erkennung. \\
\texttt{isls-reader/src/go.rs}
  & Go-Parser: \texttt{parse\_go()}, Receiver-Methods, 
    Import-Blocks, Sichtbarkeit via Gro{\ss}buchstabe. \\
\texttt{isls-reader/src/normalize.rs}
  & \texttt{normalize\_name()}, Return-Type-Unwrapping, 
    gemeinsame Normalisierungslogik. \\
\bottomrule
\end{longtable}


% ══════════════════════════════════════════════════════════════════════════════
\section{Modifizierte Dateien}
% ══════════════════════════════════════════════════════════════════════════════

\begin{longtable}{lp{10cm}}
\toprule
\textbf{File} & \textbf{Change} \\
\midrule
\texttt{isls-reader/src/lib.rs}
  & \texttt{pub mod python; pub mod typescript; pub mod go; 
    pub mod normalize;} Erweitere \texttt{Language}-Enum. 
    \"Andere \texttt{parse\_string()} und 
    \texttt{parse\_directory()} f\"ur Multi-Language. \\
\texttt{isls-code-topo/src/lib.rs}
  & \texttt{CodeTopology} bekommt \texttt{languages} 
    und \texttt{primary\_language}. \\
\texttt{isls-norms/src/learning.rs}
  & \texttt{NormCandidate} bekommt \texttt{languages} 
    und \texttt{cross\_language}. 
    Promotion-Check: wenn cross-language, logge es. \\
\texttt{isls-gateway/src/discover.rs}
  & GitHub-Suche: Sprachfilter entfernen oder erweitern. 
    X-Ray: \texttt{languages}-Feld anzeigen. \\
\texttt{isls-forge-llm/src/codematrix.rs}
  & \texttt{extract\_resonites()}: Language-Parameter 
    durchreichen. \\
\bottomrule
\end{longtable}


% ══════════════════════════════════════════════════════════════════════════════
\section{Constraints for Claude Code}
% ══════════════════════════════════════════════════════════════════════════════

\begin{enumerate}
\item \textbf{Regex-basierte Parser.} Keine externen 
      AST-Bibliotheken (kein tree-sitter, kein syn f\"ur 
      Python). Reines Regex, Zeile f\"ur Zeile. Das ist 
      bewusst: wir brauchen nur Top-Level-Deklarationen, 
      keinen vollst\"andigen AST.
\item \textbf{Alle Parser produzieren identische 
      \texttt{Resonite}-Varianten.} Keine neuen Varianten 
      f\"ur Python/TS/Go. Die bestehenden 5 Varianten 
      (Fn, Type, Import, Layer, Relation) reichen.
\item \textbf{Name-Normalisierung ist zentral.} Alle Parser 
      rufen \texttt{normalize\_name()} auf. camelCase, 
      PascalCase, snake\_case werden vereinheitlicht.
\item \textbf{Return-Type-Unwrapping} entfernt 
      Sprach-Wrapper: \texttt{Result<T>} $\to$ \texttt{T}, 
      \texttt{Promise<T>} $\to$ \texttt{T}, 
      \texttt{Optional[T]} $\to$ \texttt{T}, 
      \texttt{(*T, error)} $\to$ \texttt{T}.
\item \textbf{Der bestehende Rust-Parser wird nicht 
      ge\"andert.} Nur das Routing in \texttt{parse\_string()} 
      wird erweitert.
\item \textbf{Vendor-Verzeichnisse} (\texttt{node\_modules}, 
      \texttt{\_\_pycache\_\_}, \texttt{venv}, \texttt{vendor}, 
      etc.) werden beim Directory-Walk ignoriert.
\item \textbf{Do not modify} isls-merkaba, isls-gateway, 
      isls-chat, isls-hypercube, isls-types, isls-cli, 
      isls-forge-llm (au{\ss}er codematrix.rs Language-Parameter).
\item All existing tests MUST pass.
\item Branch: \texttt{claude/implement-isls-i4-spec}.
\item Commit format: \texttt{I4/W\{n\}: <description>}.
\end{enumerate}


% ══════════════════════════════════════════════════════════════════════════════
\section{Tests}
% ══════════════════════════════════════════════════════════════════════════════

Pro Parser mindestens 5 Unit-Tests:

\begin{enumerate}[nosep]
\item Parse eine Datei mit 3+ Funktionen $\to$ korrekte 
      Fn-Resonites.
\item Parse eine Datei mit Classes/Structs $\to$ korrekte 
      Type-Resonites mit field\_count.
\item Parse Imports $\to$ korrekte Import-Resonites mit 
      is\_external.
\item Async-Erkennung $\to$ is\_async korrekt.
\item Sichtbarkeit $\to$ is\_pub korrekt.
\end{enumerate}

Plus Normalisierungs-Tests:

\begin{enumerate}[nosep]
\item \texttt{camelCase} $\to$ \texttt{camel\_case}.
\item \texttt{PascalCase} $\to$ \texttt{pascal\_case}.
\item \texttt{snake\_case} $\to$ \texttt{snake\_case} (unver\"andert).
\item \texttt{Result<User>} $\to$ \texttt{User}.
\item \texttt{Promise<User>} $\to$ \texttt{User}.
\end{enumerate}

Plus ein Integration-Test:

\begin{enumerate}[nosep]
\item \texttt{parse\_directory()} auf ein Testverzeichnis mit 
      je einer .rs, .py, .ts, .go Datei $\to$ alle vier 
      Sprachen erkannt, Resonites aus allen vier extrahiert.
\end{enumerate}


% ══════════════════════════════════════════════════════════════════════════════
\section{Acceptance Criteria}
% ══════════════════════════════════════════════════════════════════════════════

\begin{longtable}{clp{9cm}}
\toprule
\textbf{\#} & \textbf{Pass} & \textbf{Criterion} \\
\midrule
1 & & \texttt{cargo build --workspace} compiles. \\
2 & & \texttt{cargo test --workspace} passes (inkl. 15+ neue Tests). \\
3 & & Python-Parser: 3 Funktionen + 2 Klassen + 3 Imports korrekt. \\
4 & & TypeScript-Parser: Arrow-Functions + Interfaces erkannt. \\
5 & & Go-Parser: Exported vs. unexported korrekt (Gro{\ss}buchstabe). \\
6 & & \texttt{parse\_directory()} erkennt alle 4 Sprachen. \\
7 & & Normalisierung: camelCase/PascalCase $\to$ snake\_case. \\
8 & & Scraping eines Python-Repos produziert Resonites + Observations. \\
9 & & Cross-language Candidate hat \texttt{cross\_language: true}. \\
10 & & Alle bestehenden Tests bestehen (backward compat). \\
\bottomrule
\end{longtable}


\vfill
\begin{center}
\color{isls-gray}\rule{0.5\textwidth}{0.4pt}\\[0.5em]
{\small End of I4 Specification.\\[0.3em]
Eine Funktion in Rust.\\
Dieselbe Funktion in Python.\\
Dieselbe Funktion in Go.\\
Drei Sprachen. Ein Resonite. Eine Norm.\\
Die Struktur ist sprachunabh\"angig.\\
Die Erkenntnis auch.}
\end{center}

\end{document}
